\documentclass[10pt]{article}

\usepackage[margin=1.8cm]{geometry}
\usepackage{enumitem}
\usepackage{titlesec}

\setlist[itemize]{
    itemsep=6pt,
    topsep=4pt,
    parsep=0pt,
    partopsep=0pt,
    leftmargin=1.4em
}

\titlespacing*{\section}{0pt}{10pt}{6pt}
\titlespacing*{\subsection}{0pt}{8pt}{4pt}
\titlespacing*{\paragraph}{0pt}{4pt}{4pt}

\titleformat{\section}
  {\normalfont\large\bfseries}
  {}
  {0pt}
  {\titlerule\vspace{0.5ex}}

\setlength{\parindent}{0pt}
\setlength{\parskip}{4pt}

\newcommand{\mainsection}[1]{%
  \begin{center}
    \Large\bfseries #1
  \end{center}
  \vspace{6pt}
}
\newcommand{\tlabel}[1]{{\large\bfseries #1}}

\title{IHTC Model Transformations}
\author{Maksym Byelko}
\date{June 2026}

\begin{document}

\maketitle
\mainsection{MOVES / MODEL TRANSFORMATIONS}

\section*{Decisions}
\paragraph{D1:} What day in the planning horizon will each patient be admitted?
\paragraph{D2:} What operating theatre will each patient use?
\paragraph{D3:} What recovery room will each patient use?
\paragraph{D4:} What rooms will each nurse cover on their assigned shifts?

\section*{Model Transformations}

\subsection{Patient Admission and Removal}
\begin{itemize}
    \item \tlabel{T1.}\\
    \textbf{Type:} atomic.\\
    \textbf{Description:} Insert a random non-mandatory patient with a random admission day, room, and theatre.\\
    \textbf{Classes used:} \texttt{Patient}, \texttt{Admission}, \texttt{Room}, \texttt{OperatingTheatre}.\\
    \textbf{Purpose:} to increase the number of scheduled patients and explore new possible allocations quickly.
    
    \item \tlabel{T2.}\\
    \textbf{Type:} atomic.\\
    \textbf{Description:} Insert the non-mandatory patient with the least workload into a feasible admission day, room, and theatre.\\
    \textbf{Classes used:} \texttt{Patient}, \texttt{Admission}, \texttt{Room}, \texttt{OperatingTheatre}, \texttt{SurgeonAvailability}, \texttt{OperatingTheatreAvailability}, \texttt{RoomAvailability}.\\
    \textbf{Purpose:} to add a patient who is easier to fit into the schedule.
    
    \item \tlabel{T3.}\\
    \textbf{Type:} atomic.\\
    \textbf{Description:} Find the earliest available surgeon, then find a theatre and room available on the same day, and insert a random compatible non-mandatory patient.\\
    \textbf{Classes used:} \texttt{Patient}, \texttt{Admission}, \texttt{Surgeon}, \texttt{SurgeonAvailability}, \texttt{OperatingTheatre}, \texttt{OperatingTheatreAvailability}, \texttt{Room}, \texttt{RoomAvailability}.\\
    \textbf{Purpose:} to insert a patient in a way that is more strongly guided by actual resource availability, especially surgeon and theatre constraints.
    
    \item \tlabel{T4.}\\
    \textbf{Type:} atomic.\\
    \textbf{Description:} Remove a random non-mandatory patient.\\
    \textbf{Classes used:} \texttt{Patient}, \texttt{Admission}.\\
    \textbf{Purpose:} to free up resources while only sacrificing an optional patient, which helps repair overloaded or poor-quality schedules.
    
    \item \tlabel{T5.}\\
    \textbf{Type:} atomic.\\
    \textbf{Description:} Remove a random patient.\\
    \textbf{Classes used:} \texttt{Patient}, \texttt{Admission}.\\
    \textbf{Purpose:} to make a larger change to the schedule and free capacity, even if that means removing a more important patient.
    
    \item \tlabel{T6.}\\
    \textbf{Type:} joint (T4 + T1).\\
    \textbf{Description:} Remove a random non-mandatory patient, then insert a random non-mandatory patient.\\
    \textbf{Classes used:} \texttt{Patient}, \texttt{Admission}, \texttt{Room}, \texttt{OperatingTheatre}, \texttt{SurgeonAvailability}, \texttt{OperatingTheatreAvailability}, \texttt{RoomAvailability}.\\
    \textbf{Purpose:} to swap one optional patient for another, testing whether a different patient gives a better schedule without changing the total number of optional patients.
    
    \item \tlabel{T7.}\\
    \textbf{Type:} joint (T5 + T1).\\
    \textbf{Description:} Remove a random patient, then insert a random non-mandatory patient.\\
    \textbf{Classes used:} \texttt{Patient}, \texttt{Admission}, \texttt{Room}, \texttt{OperatingTheatre}, \texttt{SurgeonAvailability}, \texttt{OperatingTheatreAvailability}, \texttt{RoomAvailability}.\\
    \textbf{Purpose:} to replace part of the current schedule with a new optional patient, allowing a more disruptive change in the set of admitted patients.
\end{itemize}

\subsection{Patient Reassignment}
\begin{itemize}
    \item \tlabel{T8.}\\
    \textbf{Type:} atomic.\\
    \textbf{Description:} Give a random admitted patient a random new compatible room.\\
    \textbf{Classes used:} \texttt{Patient}, \texttt{Admission}, \texttt{Room}, \texttt{RoomAvailability}.\\
    \textbf{Purpose:} to reassign a patient to a different feasible room in order to improve room usage or reduce incompatibilities.
    
    \item \tlabel{T9.}\\
    \textbf{Type:} atomic.\\
    \textbf{Description:} Give a random admitted patient a random new compatible admission day.\\
    \textbf{Classes used:} \texttt{Patient}, \texttt{Admission}, \texttt{SurgeonAvailability}, \texttt{OperatingTheatreAvailability}, \texttt{RoomAvailability}.\\
    \textbf{Purpose:} to move a patient to a different feasible day in order to improve schedule balance or resource usage.
    
    \item \tlabel{T10.}\\
    \textbf{Type:} atomic.\\
    \textbf{Description:} Give a random admitted patient a random new compatible theatre.\\
    \textbf{Classes used:} \texttt{Patient}, \texttt{Admission}, \texttt{OperatingTheatre}, \texttt{OperatingTheatreAvailability}.\\
    \textbf{Purpose:} to reassign a patient to a different feasible theatre in order to improve theatre allocation.
    
    \item \tlabel{T11.}\\
    \textbf{Type:} joint (T8 + T9 + T10).\\
    \textbf{Description:} Change a patient's room, admission day, and theatre sequentially.\\
    \textbf{Classes used:} \texttt{Patient}, \texttt{Admission}, \texttt{Room}, \texttt{RoomAvailability}, \texttt{OperatingTheatre}, \texttt{OperatingTheatreAvailability}, \texttt{SurgeonAvailability}.\\
    \textbf{Purpose:} to make a broader coordinated change to one patient’s allocation across several decision variables.
    
    \item \tlabel{T12.}\\
    \textbf{Type:} joint (T8 + T9).\\
    \textbf{Description:} Change a patient's room and admission day sequentially.\\
    \textbf{Classes used:} \texttt{Patient}, \texttt{Admission}, \texttt{Room}, \texttt{RoomAvailability}, \texttt{SurgeonAvailability}, \texttt{OperatingTheatreAvailability}.\\
    \textbf{Purpose:} to jointly adjust accommodation and timing in a way that may improve feasibility or quality more than changing one variable alone.
    
    \item \tlabel{T13.}\\
    \textbf{Type:} joint (T9 + T10).\\
    \textbf{Description:} Change a patient's admission day and theatre sequentially.\\
    \textbf{Classes used:} \texttt{Patient}, \texttt{Admission}, \texttt{OperatingTheatre}, \texttt{OperatingTheatreAvailability}, \texttt{SurgeonAvailability}, \texttt{RoomAvailability}.\\
    \textbf{Purpose:} to jointly modify the timing and surgical allocation of a patient while keeping the room unchanged.
\end{itemize}

\subsection{Nurse Assignment}
\begin{itemize}
    \item \tlabel{T14.}\\
    \textbf{Type:} atomic.\\
    \textbf{Description:} Add a random nurse to a random compatible room on a random compatible shift.\\
    \textbf{Classes used:} \texttt{Nurse}, \texttt{NurseWorkingShift}, \texttt{HospitalisationShift}, \texttt{Room}, \texttt{RoomShiftAssignment}.\\
    \textbf{Purpose:} to create a new nurse-to-room assignment where the nurse is available on that shift and satisfies the required skill and workload conditions.
    
    \item \tlabel{T15.}\\
    \textbf{Type:} atomic.\\
    \textbf{Description:} Remove a random nurse from a random room on their schedule.\\
    \textbf{Classes used:} \texttt{Nurse}, \texttt{HospitalisationShift}, \texttt{RoomShiftAssignment}.\\
    \textbf{Purpose:} to remove an existing nurse-to-room assignment and allow the optimisation to repair or rebalance staffing decisions.
    
    \item \tlabel{T16.}\\
    \textbf{Type:} joint (T15 + T14).\\
    \textbf{Description:} Remove a nurse from one room and add them to another.\\
    \textbf{Classes used:} \texttt{Nurse}, \texttt{NurseWorkingShift}, \texttt{HospitalisationShift}, \texttt{Room}, \texttt{RoomShiftAssignment}.\\
    \textbf{Purpose:} to reallocate nursing capacity between room-shifts without changing the total number of nurse assignments.
\end{itemize}

\end{document}
